\section{Paraxial Refractive Solution}
\label{sec:paraxial}

We now seek the refractive-index profile that reproduces the leading-order
horizon geometry within the paraxial approximation.

\subsection{Small-angle ray equation}

We first restrict attention to rays with small slope relative to the
horizontal, as in standard paraxial treatments of stratified atmospheric
refraction \cite{nener2003analytical},
\begin{equation}
    |h'|\ll 1.
\end{equation}
Expanding the Fermat integrand in Eq.~\eqref{eq:fermat-functional} gives
\begin{equation}
    n(h)\sqrt{1+h'^2}
    =
    n(h)\left(1+\frac{h'^2}{2}\right)
    +\mathcal{O}(h'^4).
    \label{eq:paraxial-lagrangian}
\end{equation}
The Euler--Lagrange equation for the retained terms in
Eq.~\eqref{eq:paraxial-lagrangian} is
\begin{equation}
    \frac{d}{dx}\bigl(n(h)h'\bigr)
    -\frac{dn}{dh}\left(1+\frac{h'^2}{2}\right)
    =0.
\end{equation}
Neglecting terms of order $h'^2$ then yields
\begin{equation}
    n(h)h''\simeq \frac{dn}{dh},
\end{equation}
or, equivalently,
\begin{equation}
    h''\simeq \frac{d\ln n}{dh}.
    \label{eq:paraxial-ray-equation}
\end{equation}
Thus, within the paraxial approximation, the logarithmic vertical gradient
of the refractive index determines the local ray curvature.

\subsection{Profile producing constant ray curvature}

Relative to a local horizontal tangent plane, the leading-order horizon
geometry of a sphere of radius $R$ is generated by curves with constant upward
curvature $1/R$. To reproduce that ray family above a flat
boundary, we therefore impose
\begin{equation}
    h''=\kappa,
    \qquad
    \kappa>0,
    \label{eq:constant-curvature-condition}
\end{equation}
with $\kappa=1/R$ for the spherical comparison. Substituting Eq.~\eqref{eq:constant-curvature-condition} into
Eq.~\eqref{eq:paraxial-ray-equation} gives
\begin{equation}
    \frac{d\ln n}{dh}=\kappa.
\end{equation}
Integrating with the boundary condition $n(0)=n_0$ yields
\begin{equation}
    n(h)=n_0 e^{\kappa h}.
    \label{eq:paraxial-exponential-profile}
\end{equation}
For the spherical comparison scale,
\begin{equation}
    \kappa=\frac{1}{R},
\end{equation}
so that
\begin{equation}
    n(h)=n_0 e^{h/R}.
    \label{eq:earth-scale-exponential-profile}
\end{equation}
The exponential form is therefore derived from the constant-curvature
requirement rather than assumed. The required refractive index increases with
height, a sign that will be central to the later comparison with an ordinary
atmosphere.

\subsection{Parabolic rays and the visibility law}

The corresponding ray family follows by integrating
Eq.~\eqref{eq:constant-curvature-condition}:
\begin{equation}
    h(x)=h_m+\frac{\kappa}{2}(x-x_m)^2,
    \label{eq:paraxial-ray-family}
\end{equation}
where $(x_m,h_m)$ is the vertex of the parabola. A limiting ray that just
grazes the opaque boundary has $h_m=0$, and its tangency point may be chosen
as $x_m=0$. Its two branches are therefore
\begin{equation}
    h(x)=\frac{\kappa x^2}{2}.
    \label{eq:paraxial-grazing-ray}
\end{equation}
An endpoint at height $h$ lies on this grazing ray at horizontal distance
\begin{equation}
    d_{\mathrm{p}}(h)
    =
    \sqrt{\frac{2h}{\kappa}}.
    \label{eq:paraxial-horizon-distance}
\end{equation}
Consequently, for observer and target heights $h_{\mathrm{o}}$ and
$h_{\mathrm{t}}$, the maximum separation is
\begin{equation}
    L_{\max}^{(\mathrm{p})}
    =
    \sqrt{\frac{2h_{\mathrm{o}}}{\kappa}}
    +
    \sqrt{\frac{2h_{\mathrm{t}}}{\kappa}}.
    \label{eq:paraxial-maximum-separation}
\end{equation}
Setting $\kappa=1/R$ gives
\begin{equation}
    L_{\max}^{(\mathrm{p})}
    =
    \sqrt{2Rh_{\mathrm{o}}}
    +
    \sqrt{2Rh_{\mathrm{t}}}.
    \label{eq:paraxial-spherical-law}
\end{equation}
This is precisely the leading-order spherical horizon law for heights much
smaller than $R$.

The paraxial construction therefore establishes that a flat-boundary model
can reproduce the leading spherical horizon-occlusion law: the lower
parts of a distant object become inaccessible first, and raising either
endpoint increases the limiting visible distance according to the same
square-root height dependence as on a sphere. The agreement is asymptotic,
however. The next section solves the exponential profile exactly
and determines how far the agreement extends beyond the small-angle regime.
